\section{Survey Performance}
\label{sec:performance}

\subsection{Sky coverage and cadence}
\label{sec:coverage}

The first-year observations contain \yrNexp\ exposures on
\yrNnight\ nights, with a total open-shutter time of \yrHours~hr.
The latter is the sum of exposure durations and describes the time
spent collecting light. Table~\ref{tab:results} summarizes the main
measurements of the survey.

We calculate each exposure's footprint from the image WCS, including
its TPV distortion terms, and combine overlapping polygons on a
HEALPix grid \citep{Gorski2005}. At $N_{\rm side}=1024$, the area
observed at least once is $\yrArea~\degsq$. Counting exposures
within each observing night and then combining sky pixels that
receive at least three visits gives $\yrAreaThree~\degsq$.
These areas describe the geometric footprint. Figure~\ref{fig:coverage}
compares the exposure coverage with the sky distribution of the
known-asteroid measurements. The exposure map traces the observing
pattern, while the asteroid map also reflects the population,
viewing geometry, and brightness selection.

\begin{figure}[tbp]
\centering
\annualplot{fig03_coverage_and_known_sky.pdf}{\linewidth}
\caption{First-year observations in J2000 ecliptic longitude and
latitude. (a) Number of science exposures covering each HEALPix pixel,
using image WCS footprints and $N_{\rm side}=512$.
(b) Number of primary $1\asec$ known-asteroid associations per pixel
at $N_{\rm side}=128$. The orange curve marks the celestial equator.
Both color scales are logarithmic; their counts refer to the different
quantities and pixel areas indicated. Geometric areas in the text
use $N_{\rm side}=1024$.}
\label{fig:coverage}
\end{figure}

Within a night, exposures are grouped by field identifier and actual
pointing. Every pair of centers in a group must be separated by at
most $0.30^\circ$, approximately one tenth of the camera width.
This separates observations with substantially different pointings
while preserving repeated visits to the same field.

The \yrNGroups\ field--night groups comprise \yrNOne\ single visits,
\yrNTwo\ pairs, and \yrNThree\ sequences with at least three visits.
Thus \yrFracThree\% of observed field--night groups provide three
or more measurements. The median interval between adjacent visits
is \yrCadence~min, with a 16th--84th percentile range of
\yrCadenceLow--\yrCadenceHigh~min. Sequences with at least three
visits have a median first-to-last span of \yrSpan~min and a
corresponding range of \yrSpanLow--\yrSpanHigh~min.

Figure~\ref{fig:cadence} shows how the observations are distributed
through the year and within individual nights. The fraction of
complete field--night sequences and the area ever observed three
times in a night describe complementary aspects of the sampling.
The former measures how often an observed field supplies a sequence
suitable for motion linking; the latter measures the accumulated
sky area with at least one such opportunity.

\begin{figure}[tbp]
\centering
\annualplot{fig04_activity_and_cadence.pdf}{0.96\linewidth}
\caption{Observing activity and cadence during the first year.
(a,b) Nightly and monthly numbers of standard survey exposures
and those entering the science sample. All panels use the same
local-noon observing-night definition. (c) Number of visits in
each field--night pointing group. (d) Exposure-midpoint separation
of adjacent visits within those groups. (e) First-to-last midpoint
span for groups with at least three visits. The time axis covers
the full first-year interval. The interval panels emphasize the
main distributions.}
\label{fig:cadence}
\end{figure}

\begin{table}[tbp]
\centering
\caption{First-Year Survey Measurements}
\label{tab:results}
\begin{tabular}{ll}
\toprule
Quantity & Measured value\\
\midrule
Science exposures / observing nights & \yrNexp\ / \yrNnight\\
Open-shutter time & \yrHours~hr\\
Geometric area observed at least once & $\yrArea~\degsq$\\
Area with three visits in at least one night & $\yrAreaThree~\degsq$\\
Groups with at least three visits & \yrNThree\ / \yrNGroups\ (\yrFracThree\%)\\
Median adjacent-visit interval & \yrCadence~min\\
Median span of sequences with at least three visits & \yrSpan~min\\
Median catalog $5\sigma$ model depth & \yrDepth~mag\\
Model-depth interquartile range & \yrDepthLow--\yrDepthHigh~mag\\
\midrule
Primary associations / catalog identities & \yrNAssoc\ / \yrNKnown\\
Identities in the main-belt orbital region & \yrNMainBelt\ (\yrMainBeltFrac\%)\\
Median radial O$-$C & $\yrOC\asec$\\
Robust scatter in $\Delta\alpha\cos\delta$ / $\Delta\delta$ & $\yrRAScatter\asec$ / $\yrDecScatter\asec$\\
\midrule
Linked / eligible known-object nights & \yrNLinked\ / \yrNEligible\ (\yrRetention\%)\\
Image-reviewed single-night candidates & \yrNCandidates\\
Measurements / nights in retained candidates & \yrNCandidateMeasurements\ / \yrNCandidateNights\\
\bottomrule
\end{tabular}
\tablecomments{Depth is defined by Equation~\ref{eq:depth}. Astrometric
statistics refer to the $1\asec$ association sample; robust scatter
is the Gaussian-equivalent interquartile spread. The linkage counts
use the observations and eligibility criteria in
Section~\ref{sec:retention_method}.}
\end{table}

\subsection{Photometric and astrometric precision}
\label{sec:photometry}

The images have a median FWHM of $\yrFWHM\asec$, with an
interquartile range of $\yrFWHMLow$--$\yrFWHMHigh\asec$.
We characterize the per-exposure sensitivity using the catalog's
background-noise model,
\begin{align}
F_{5,\rm ADU} &=\frac{5}{g_{\rm e}}
 \left[(s_{\rm ADU}g_{\rm e}+r_{\rm e}^{\,2})\,
 \pi\left(\frac{w_{\rm pix}}{2}\right)^2\right]^{1/2},\\
m_{5,\rm model} &= ZP-2.5\log_{10}F_{5,\rm ADU},
\label{eq:depth}
\end{align}
where $s_{\rm ADU}$ is the sky level per pixel, $g_{\rm e}$ the gain
in electrons per ADU, $r_{\rm e}$ the read noise in electrons,
and $w_{\rm pix}$ the image FWHM in pixels. The noise area is a
circle of radius FWHM/2, and $ZP$ is the zero point from the
18-pixel-diameter calibration aperture. The calculation includes
sky and read noise, with no additional correction for the
encircled energy at FWHM/2. Its median value is \yrDepth~mag,
with an interquartile range of \yrDepthLow--\yrDepthHigh~mag
(Figure~\ref{fig:quality}a). This model depth describes the image
background sensitivity; the fixed eight-pixel measurement aperture
has its own flux uncertainty.

Figure~\ref{fig:quality}b shows the aperture-magnitude uncertainties
for unflagged known-asteroid associations with positive measured
fluxes and uncertainties. Selected stellar sequences provide a
separate check on photometric repeatability. We use 18 images in
six field--night sequences and compare isolated, unflagged stars
between visits. The single-epoch scatter is estimated as
\begin{equation}
\sigma_{\rm epoch}=\frac{\mathrm{IQR}(\Delta g_{\rm aper})}
 {2\Phi^{-1}(0.75)\sqrt{2}},
\label{eq:repeatability}
\end{equation}
where $\Phi^{-1}$ is the standard-normal quantile function.
This expresses the interquartile spread as a Gaussian-equivalent
scatter and assigns comparable noise to the two epochs.

For $15.5\leq\gaper<16$, the scatter is 11.77~mmag, compared
with a median reported uncertainty of 11.53~mmag. At
$17.5\leq\gaper<18$, the corresponding values are
56.89 and 57.41~mmag. These selected sequences provide a
within-night repeatability check under relatively clean imaging
conditions. The first-year distribution of catalog uncertainties
in Figure~\ref{fig:quality}b describes the wider range of
observing conditions.

\label{sec:astrometry}
We compare source centroids with the predicted asteroid positions
at each exposure midpoint. The right-ascension difference is
wrapped to $[-180^\circ,180^\circ]$ and multiplied by the cosine
of the predicted declination. The \yrNAssoc\ primary associations
have a median radial O$-$C of $\yrOC\asec$. Their signed component
medians are \yrRAOffset~mas in $\Delta\alpha\cos\delta$ and
\yrDecOffset~mas in $\Delta\delta$, with Gaussian-equivalent
interquartile scatters of $\yrRAScatter\asec$ and
$\yrDecScatter\asec$, respectively.

The magnitude dependence of the radial separation is shown in
Figure~\ref{fig:quality}c. Panel~(d) shows the nightly medians of
the signed components and the number of measurements in each
night. These comparisons probe the combined effects of centroid
measurement, the astrometric reference frame, observation time,
and the predicted position. The quoted statistics describe
associations selected within $1\asec$.

\begin{figure}[tbp]
\centering
\annualplot{fig05_measurement_quality.pdf}{\linewidth}
\caption{Photometric and astrometric measurements during the first
year. (a) Catalog $5\sigma$ model depth in 0.2-mag bins; the orange
line and band mark the median and interquartile range.
(b) Aperture-magnitude uncertainty for unflagged primary associations
with positive magnitude uncertainty, flux, and flux uncertainty.
Cells are 0.1~mag by 0.04~dex. Orange curves show the median and
16th--84th percentiles in 0.5-mag bins with at least 30 measurements.
Blue points show the stellar repeatability from the six selected
sequences, using bins with at least 50 stellar pairs.
(c) Radial O$-$C in 0.1-mag by $0.02\asec$ cells, with the same
trend conventions. The color bar gives the number of measurements
per cell. These associations lie within $1\asec$ of the predictions.
(d) Nightly medians of $\Delta\alpha\cos\delta$ and $\Delta\delta$;
the gray bars show the number of associations per night.
The magnitude-dependent panels use measurements with usable
\gaper\ in the displayed range.}
\label{fig:quality}
\end{figure}

\subsection{Linkage of repeatedly detected asteroids}
\label{sec:retention_results}

The first-year assessment contains \yrNEligible\ eligible known-object
nights, of which \yrNLinked\ are represented by motion-search links.
The conditional retention is therefore $R_{\rm cond}=\yrRetention\%$.
The 95\% Wilson interval is \yrRetentionLow--\yrRetentionHigh\%;
it describes binomial counting uncertainty for the evaluated sample.
The magnitude and angular-rate dependence is shown in
Figure~\ref{fig:retention}.

The two-dimensional count map shows where the detections constrain
the measurement in median object--night \gaper\ and predicted
angular rate. The adjacent map gives the linked fraction in the
same cells. The lower panels summarize the dependence on each
quantity and show the corresponding eligible counts. These
distributions relate the observed linking behavior to the brightness,
displacement, and temporal sampling available to the search.

\begin{figure}[tbp]
\centering
\annualplot{fig06_conditional_retention.pdf}{\linewidth}
\caption{Conditional linkage retention for first-year known-object
nights. Eligibility requires detections in at least three distinct
science exposures within $1.5\asec$ of one known identity, with
predicted first-to-last angular rate $3$--$63$~\rateunit.
(a) Eligible counts in 0.25-mag by 3~\rateunit\ cells, with a
logarithmic color scale. (b) Linked fractions in the same cells;
cells containing fewer than 30 eligible object--nights are gray.
(c,d) One-dimensional linked fractions with 95\% Wilson intervals
and horizontal bars indicating bin widths. The lower histograms
show eligible counts per bin. The dashed line marks the overall
fraction, \yrRetention\%. Objects with unusable magnitudes remain
in the overall and angular-rate samples and are omitted only from
panels requiring a magnitude. Both maps use unsmoothed counts.}
\label{fig:retention}
\end{figure}
